\runninghead{Scope and References}
\section{Appendix: Scope and Qualifications}
\label{sec:scope}

{\small
\textbf{Edition, Mass and occurrence.} The study treats the Mass
\latin{Dominica vigesima quarta et ultima post Pentecosten}, II classis, of
the \work{Missale Romanum}, editio typica, Typis Polyglottis Vaticanis 1962,
printed pp.~422--424, nos.~1769--1778, as said on Sunday 22~November 2026 in
the universal 1962 calendar and no particular calendar. Its five chants are
printed in the same words as those of the Twenty-third Sunday after Pentecost
(pp.~416--417); its orations and lessons are its own. Its place in 2026, the
omission of the feast of St Cecilia, the colour, the Gloria, the Credo and the
Trinity Preface rest on the Missal's general rubrics (RG~10, 12, 15--16, 18~b,
91, 107--111, 127~b; RGMR~434, 494~b), on the note closing the Twenty-third
Sunday (p.~417) and on the Missal's printed direction after the Secret, as
\texttt{research/context.md} records. Two dated Ordos of institutes using the
1962 books, of the Fraternity of St Peter and of the Institute of Christ the
King Sovereign Priest, both in France, give the same Mass for the date, of the
second class and in green. No local oration, external solemnity or votive Mass
is considered. The \latin{Asperges}, the Ordinary and the Office are not
treated.

\textbf{Text control.} Every Latin form was read on page images of the CMAA
facsimile of the 1962 typical edition and collated word for word with the
Ratisbon printing of 1862 (Pustet), whose page images carry every word of the
ten texts; the collation and the comparison with the Clementine Vulgate are in
\texttt{propers/verified.md}. Where the chants and lessons differ from the
Clementine Vulgate the differences are reported as textual facts; the Latin
psalter, antiphonary or lectionary from which those forms came was not
identified. The history of this Gospel's assignment to the last Sunday after
Pentecost is not treated. The rubrics of the 1962 book are described and not
reproduced.

\textbf{English.} The English of the two lessons, of the Introit's psalm
verse and of each chant's Bible verses is the Douay--Rheims in Challoner's
revision. The English of the five chants as the Missal sings them and of the
three orations is that of the Lasance \work{New Roman Missal} in the Benziger
revision of 1945, public domain in the United States by non-renewal. None of
this English is an approved liturgical translation of the 1962 Missal.
Chrysostom's homilies on Matthew and Colossians are quoted in the English of
the Nicene and Post-Nicene Fathers; the Greek was not read. Augustine on
Psalm 129 is quoted in the same series' English, which is abridged, so that
those quotations are of the condensed English and not of Augustine's
continuous text. Dom Lucien Fromage is quoted in the English of Dom Laurence
Shepherd. Elsewhere the Latin of the Fathers and medieval writers is quoted
with an English rendering made for this study. Where the Lasance English
departs from the Latin of a chant or an oration, single Latin words or short
phrases are glossed literally; no English version of a chant or an oration is
supplied.

\textbf{Reception.} The interpretations rest on the loci named in the text and
recorded, with editions and bounds, in \texttt{research/scope.md} and
\texttt{research/interpretations.md}. Every connection drawn between two texts
of the Mass belongs to the study's interpretations unless an author is named as
making it; no Father cited here comments on this Mass as a whole, and none of
them speaks of the Eucharist at Mt 24:28. Gregory is cited on Luke's form of
the Gospel's discourse (Lk 21:25--33) and Bede on Mark's (Mk 13:14--31); where
Bede repeats Jerome word for word, Jerome is cited. The Latin of Jerome on
Matthew and on Jeremias was read on page images; Augustine, Gregory and Bede in
transcriptions of Migne not collated with page images; Hilary, Aquinas on
Matthew, Rupert and Sicard in uncorrected optical text, whose quotations here
are normalised in spelling; Aquinas on Colossians and on Psalm 43 in the Corpus
Thomisticum text. Every principal author of an interpretation has a row in the
repository's register of authorial standing. Rupert of Deutz, Sicard of
Cremona and Lucien Fromage are ecclesiastical writers and are cited in support
only. Rupert and Sicard expound these chants in Masses whose lessons differed
from this Mass's, and they are cited for the chants alone; William Durandus's
account of the same chants repeats Sicard's and is not cited. No commentary on
the Collect, Secret or Postcommunion survives among the commentators read
beyond Fromage's brief notes. Origen on Matthew, the psalm commentaries of
Cassiodorus and Bellarmine on these chants, and the Greek commentators on
Colossians were not read.

\textbf{Chronology.} Every biblical date in the Scriptural Date and Location
appendix is generated from the repository's Scripture chronology corpus,
with its sources named there.

\textbf{Rights.} The Latin is published on 17 U.S.C. 103(b) together with its
public-domain witness, the Pustet printing of 1862. The Douay--Rheims and the
Lasance English are quoted from registered editions. The other quotations are
short and come from editions published before 1931 or from medieval and
ancient Latin: Migne's \work{Patrologia Latina}, vols.~9, 24, 26, 92, 170 and
213, and the Migne texts of Augustine and Gregory; the Venice Aquinas (1745);
the Nicene and Post-Nicene Fathers; and the Dublin \work{Liturgical Year}
(1883). The Corpus Thomisticum and Wikisource texts are used only for Latin
that is itself medieval or ancient. Scripture, liturgical texts and quoted
translations are not relicensed by the repository's licence.

\textbf{Use.} This is a study aid to the 1962 Missal, not an official
liturgical book, translation or critical edition.
\par}

\section{References}
\label{sec:references}

{\small
\subsection*{Liturgical books}
\begin{itemize}[leftmargin=1.2em,itemsep=.25em]
\item \work{Missale Romanum ex decreto Sacrosancti Concilii Tridentini
restitutum, Summorum Pontificum cura recognitum}, editio typica (Typis
Polyglottis Vaticanis, 1962), CMAA facsimile: \work{Rubricae generales} and
\work{Rubricae generales Missalis romani}; Proprium de Tempore, pp.~416--417
and 422--424, nos.~1769--1778.
\item \work{Missale Romanum} (Ratisbon: Pustet, 1862), pp.~359--360.
\item Lasance, F.~X., and Walsh, F.~A., \work{The New Roman Missal in Latin and
English}, revised edition (New York: Benziger Brothers, 1945), pp.~736--740.
\item Fraternité Saint-Pierre (France), \work{Ordo du mois}, and Institut du
Christ Roi Souverain Prêtre (France), \work{Ordo}, entries for 22~November
2026 (websites), read 16~November 2026.
\end{itemize}

\subsection*{Scripture}
\begin{itemize}[leftmargin=1.2em,itemsep=.25em]
\item \work{The Holy Bible Translated from the Latin Vulgate} (Douay--Rheims,
Challoner revision), Project Gutenberg text: Gen 43:16, 34; 50:17--20; Ps
21:17; 23:6; 43:4--23; 75:5; 83:6; 84:1--5; 102:5; 129:1--8; Is 40:31; 49:25;
53:7; Jer 29:1--14; Dan 9:27; Hab 2:3; Mt 6:6; 7:7; 24:1--36; Mk 11:20--25;
Lk 17:10; 21:25--33; Jn 12:24; Rom 7:23; 1~Cor 15:36; 2~Cor 5:1; Col
1:3--15; Heb 10:32; 1~Jn 2:15.
\item \work{Biblia Sacra Vulgatae Editionis} (Clementine Vulgate), for the
comparisons with the Missal's forms.
\end{itemize}

\subsection*{Fathers, Doctors and other authors}
\begin{itemize}[leftmargin=1.2em,itemsep=.25em]
\item Augustine, \work{Enarrationes in Psalmos} 43.8--10 and 84.3--5, Migne
Latin text (Wikisource); 129.1--3, NPNF, first series, vol.~8 (abridged);
\work{Quaestiones Evangeliorum} I.37--39 and I.42, Migne Latin text
(Wikisource).
\item Bede, \work{In Marci evangelium expositio} III, on 11:22--25, and IV, on
13:14--16 and 13:31, Migne Latin text (Wikisource).
\item Chrysostom, John, \work{Homilies on Matthew} 75.2, 76.1--3 and 77.1,
NPNF, first series, vol.~10; \work{Homilies on Colossians} 2, NPNF, first
series, vol.~13.
\item Gregory the Great, \work{Homiliae in Evangelia} I.1, nn.~1, 3--5, Migne
Latin text (Wikisource).
\item Hilary of Poitiers, \work{In Evangelium Matthaei commentarius} 25.3--8
and 26.1--3, PL~9 (Paris, 1844), cols.~1053--1058; \work{Tractatus super
Psalmos}, on Ps 129, PL~9, col.~718.
\item Jerome, \work{Commentarii in Matthaeum} IV, on 24:15--35, PL~26 (Paris,
1845), cols.~177--181; \work{In Hieremiam} V, on 29:10--14, PL~24 (Paris,
1865), cols.~893--894.
\item Thomas Aquinas, \work{Super Evangelium S.~Matthaei lectura}, cap.~24, in
\work{Opera omnia}, vol.~3 (Venice: Bettinelli, 1745), pp.~306--310;
\work{Super epistolam ad Colossenses}, cap.~1, lect.~3, Corpus Thomisticum
text; \work{Super Psalmos}, on Ps 43, nn.~3--4, Corpus Thomisticum text.
\end{itemize}

\subsection*{Commentators on the Mass}
\begin{itemize}[leftmargin=1.2em,itemsep=.25em]
\item Fromage, Lucien, continuing Prosper Guéranger, \work{The Liturgical
Year: The Time after Pentecost}, vol.~2 (English vol.~11), trans.\ Laurence
Shepherd (Dublin: Duffy, 1883), ``The Twenty-fourth and Last Sunday after
Pentecost,'' pp.~505, 509--510, 513.
\item Rupert of Deutz, \work{De divinis officiis} XII.23, PL~170,
cols.~329--330.
\item Sicard of Cremona, \work{Mitrale} VIII.23 and VIII.25, PL~213.
\end{itemize}
\par}
